% ---
% Agradecimentos
% ---
Este trabalho não é obra individual. Devo-o a muitas pessoas e a alguns momentos que, cada um a seu modo, fizeram-no possível.

Ao Prof. Dr. Marcos Quiles, meu orientador, agradeço a generosidade acadêmica, a clareza técnica e, sobretudo, a paciência para acompanhar a evolução deste projeto — do MBA à qualificação — sem nunca abrir mão do rigor. Cada reunião foi um exercício de recalibração que este texto carrega em cada capítulo.

À minha esposa Alana Caroline, agradeço a parceria sem condições. Os fins de semana de leitura, as madrugadas de código, tudo isso só foi possível porque ela segurou, sem que eu precisasse pedir, o mundo concreto ao meu lado. E, principalmente, porque me deu a Maria Vitória.

Nossa filha nasceu em 28 de março de 2026, no meio deste mestrado. Se este trabalho tem alguma urgência ética, ela veio dessas primeiras semanas — quando eu segurava o rosto de uma pessoa novinha no mundo e pensava, sem conseguir evitar, que sistemas de reconhecimento facial precisam aprender a ver rostos como o dela com o mesmo cuidado com que eu a olhava. Este trabalho é, em parte, tentativa de dizer isso em código.

Aos meus pais, Lucia Ozzetti e Wladimir Cruz, agradeço o que não cabe em texto: a base que me permitiu chegar aqui, a curiosidade que me foi ensinada como valor e o exemplo do trabalho feito com honestidade.

Por fim, ao Programa de Pós-Graduação em Ciência da Computação da Unifesp / ICT, agradeço a estrutura acadêmica e a oportunidade de conduzir esta pesquisa com autonomia. Aos servidores da secretaria do
Programa, agradeço a paciência com os trâmites e o apoio nas tramitações administrativas.